LLM agent memories that extract facts from each conversation session often run extraction asynchronously, behind a message broker. A question can then arrive before the newest session is committed, and the reader answers from a superseded fact. Once the update is committed, stores that update in place, and read paths that mask superseded facts, lose the earlier value. We formalise both failures, read-after-write staleness and history erasure, and measure them on LongMemEval-S with a deterministic lag barrier and a commit-reordering model. We propose DeltaWAL, which writes raw turns with monotonic epochs to a write-ahead log before the broker publish and merges the uncommitted, relevance-gated log tail with stored facts in epoch order; DeltaWAL-Versioned also keeps every fact version. With one session in flight, an asynchronous store answers at most \nm{tab1/glm-5.3/B2_async/lag1/KU}\% of knowledge-update questions correctly across three readers, while DeltaWAL-Versioned answers \nm{tab1/glm-5.3/deltawal_versioned/lag1/KU}\% with glm-5.3, the same as with a complete store, and acknowledges writes in milliseconds. In-place updates and read-side masking both erase answers to questions about earlier values; \memzero{} OSS \nm{frozen_config/mem0_version} shows the same loss. On MemConflict, the system ordering replicates for genuine updates; on a small, audited set of false corrections, the newer-wins rule leads readers to believe the false statement, an exploratory limit of epoch precedence.
